% \usetikzlibrary{arrows.meta,backgrounds}
\providecommand{\chaptertitlefont}{\small\fontspec{Times New Roman}\songti}
\begin{tikzpicture}[
    x=1cm,
    y=1cm,
    font=\small\fontspec{Times New Roman}\songti,
    >=Stealth,
    line cap=round,
    line join=round,
    outerpanel/.style={
        draw=none,
        fill=gray!2,
        rounded corners=6pt,
        line width=.45pt
    },
    block/.style args={#1}{
        draw=#1!70!black,
        fill=#1!8,
        rounded corners=5pt,
        line width=.55pt
    },
    lefttitle/.style args={#1}{
        draw=none,
        fill=none,
        minimum width=2.18cm,
        minimum height=.82cm,
        align=center,
        text=black!92,
        line width=.42pt,
        font=\chaptertitlefont
    },
    chaptermaintitle/.style={
        align=center,
        text=black!92,
        font=\chaptertitlefont
    },
    card/.style={
        draw=black!26,
        fill=white,
        rounded corners=3pt,
        line width=.38pt,
        align=center,
        inner xsep=3.5pt,
        inner ysep=2.6pt,
        minimum height=.54cm,
        font=\scriptsize\fontspec{Times New Roman}\songti
    },
    bigcard/.style={
        draw=black!26,
        fill=white,
        rounded corners=3pt,
        line width=.38pt,
        align=center,
        inner xsep=4.2pt,
        inner ysep=2.8pt,
        minimum height=.56cm,
        font=\scriptsize\fontspec{Times New Roman}\songti
    },
    note/.style={
        align=center,
        text=black!72,
        font=\scriptsize\fontspec{Times New Roman}\songti,
        inner xsep=4pt,
        inner ysep=1pt,
        fill=gray!2
    },
    flow/.style={
        -{Stealth[length=3.4pt,width=3.8pt]},
        draw=black,
        line width=.58pt,
        shorten >=.4pt
    },
    auxflow/.style={
        -{Stealth[length=3.2pt,width=3.5pt]},
        draw=black,
        line width=.50pt,
        shorten >=.3pt
    },
    routebar/.style={
        draw=black,
        line width=.56pt
    },
    routeflow/.style={
        -{Stealth[length=3.4pt,width=3.8pt]},
        draw=black,
        line width=.58pt,
        shorten >=.35pt
    },
    divider/.style={draw=black!20, line width=.34pt}
]

% ------- colors -------
\colorlet{cOne}{blue}
\colorlet{cTwo}{teal}
\colorlet{cMid}{orange}
\colorlet{cFive}{purple}
\colorlet{cSix}{green!55!black}
\colorlet{cSeven}{red}

% ------- overall frame -------
\def\xmin{0.25}
\def\xmax{17.35}
\def\ymin{0.90}
\def\ymax{23.10}
\def\cx{8.80}
\def\routeL{0.34}
\def\routeR{17.26}

\begin{scope}[on background layer]
  \draw[outerpanel] (\xmin,\ymin) rectangle (\xmax,\ymax);
\end{scope}

% =====================================================
% Chapter 1
% =====================================================
\node[block=cOne, minimum width=16.45cm, minimum height=1.62cm] (bg1) at (\cx,22.00) {};
\node[chaptermaintitle] at (\cx,22.48) {第一章\quad 绪论};
\node[card, minimum width=2.45cm] (c1a) at (3.40,21.70) {研究背景\\与意义};
\node[card, minimum width=2.45cm] (c1b) at (6.95,21.70) {国内外\\研究现状};
\node[card, minimum width=2.60cm] (c1c) at (10.55,21.70) {研究问题\\与技术路线};
\node[card, minimum width=2.90cm] (c1d) at (14.20,21.70) {主要内容\\与组织结构};

% =====================================================
% Chapter 2
% =====================================================
\node[block=cTwo, minimum width=16.45cm, minimum height=1.62cm] (bg2) at (\cx,19.88) {};
\node[chaptermaintitle] at (\cx,20.36) {第二章\quad 跨工况领域自适应故障诊断基本方法};
\node[card, minimum width=1.86cm] (c2a) at (3.10,19.58) {跨工况\\任务设定};
\node[card, minimum width=1.86cm] (c2b) at (5.95,19.58) {联合分布\\最优传输};
\node[card, minimum width=1.86cm] (c2c) at (8.80,19.58) {DeepJDOT\\单源对齐};
\node[card, minimum width=1.86cm] (c2d) at (11.65,19.58) {WJDOT\\多源加权};
\node[card, minimum width=1.86cm] (c2e) at (14.50,19.58) {CoDATS\\时序表征};

\foreach \a/\b in {c2a/c2b,c2b/c2c,c2c/c2d,c2d/c2e}{
    \draw[flow] (\a.east) -- (\b.west);
}
\draw[auxflow] (bg1.south) -- (bg2.north);

% =====================================================
% Chapter 3 and 4 shared layout guide
% =====================================================
\coordinate (routeTop) at (\cx,18.82);
\node[note, anchor=west] at (9.10,18.82) {为第三至第五章提供最优传输、时序建模与领域自适应理论支撑};

% =====================================================
% Chapter 3
% =====================================================
\node[block=cMid, minimum width=16.45cm, minimum height=3.50cm] (bg3) at (\cx,16.82) {};
\node[chaptermaintitle] at (\cx,18.07) {第三章 TPU-DeepJDOT：基于时序原型非平衡联合分布对齐的单源领域自适应故障诊断方法};

\node[note] at (\cx,17.52) {时序原型非平衡联合分布对齐基本原理};
\node[card, minimum width=2.35cm] (c3a) at (4.72,16.97) {非平衡最优\\传输对齐原理};
\node[card, minimum width=2.35cm] (c3b) at (8.80,16.97) {源域类别原型\\约束原理};
\node[card, minimum width=2.35cm] (c3c) at (12.88,16.97) {时间序列统计\\代价构建原理};

\node[note] at (\cx,16.22) {基于时序原型非平衡联合分布对齐的诊断模型构建};
\node[card, minimum width=2.18cm] (c3d) at (2.50,15.62) {单源跨工况\\诊断问题建模};
\node[card, minimum width=2.18cm] (c3e) at (6.70,15.62) {非平衡联合\\传输代价构建};
\node[card, minimum width=2.18cm] (c3f) at (10.90,15.62) {源域原型与\\时序统计融合};
\node[card, minimum width=2.18cm] (c3g) at (15.10,15.62) {源域监督对比\\特征学习};

\foreach \a/\b in {c3a/c3b,c3b/c3c,c3d/c3e,c3e/c3f,c3f/c3g}{
    \draw[flow] (\a.east) -- (\b.west);
}

% =====================================================
% Chapter 4
% =====================================================
\node[block=cMid, minimum width=16.45cm, minimum height=3.50cm] (bg4) at (\cx,12.82) {};
\node[chaptermaintitle] at (\cx,14.07) {第四章 CBTPU-DeepJDOT：基于多证据置信均衡伪标签的单源领域自适应故障诊断方法};

\node[note] at (\cx,13.52) {多证据置信均衡伪标签学习基本原理};
\node[card, minimum width=2.35cm] (c4a) at (4.72,12.97) {目标域伪标签\\学习原理};
\node[card, minimum width=2.35cm] (c4b) at (8.80,12.97) {教师参考模型与\\一致性学习原理};
\node[card, minimum width=2.35cm] (c4c) at (12.88,12.97) {类别均衡\\约束原理};

\node[note] at (\cx,12.22) {基于多证据置信均衡伪标签的目标域样本利用策略};
\node[card, minimum width=2.18cm] (c4d) at (2.50,11.62) {三类语义\\分布估计};
\node[card, minimum width=2.18cm] (c4e) at (6.70,11.62) {多证据一致性\\筛选机制};
\node[card, minimum width=2.18cm] (c4f) at (10.90,11.62) {动态置信阈值\\与类别均衡};
\node[card, minimum width=2.18cm] (c4g) at (15.10,11.62) {增强一致性与\\预测融合策略};

\foreach \a/\b in {c4a/c4b,c4b/c4c,c4d/c4e,c4e/c4f,c4f/c4g}{
    \draw[flow] (\a.east) -- (\b.west);
}
\draw[routeflow] (bg3.south) -- (bg4.north);
\node[note, anchor=west] at (9.10,14.82) {单源方法递进：先建立可靠对齐，再筛选高置信目标样本};

% =====================================================
% Chapter 5
% =====================================================
\node[block=cFive, minimum width=16.45cm, minimum height=3.50cm] (bg5) at (\cx,8.815) {};
\node[chaptermaintitle] at (\cx,10.065) {第五章 CA-CCSR-WJDOT：基于类条件源域可靠性加权的多源领域自适应故障诊断方法};

\node[note] at (\cx,9.515) {类条件源域可靠性加权基本原理};
\node[card, minimum width=2.35cm] (c5a) at (4.72,8.965) {多源迁移与\\负迁移问题};
\node[card, minimum width=2.35cm] (c5b) at (8.80,8.965) {源域可靠性\\度量原理};
\node[card, minimum width=2.35cm] (c5c) at (12.88,8.965) {类条件源域\\权重分配原理};

\node[note] at (\cx,8.215) {基于类条件源域可靠性加权的多源对齐策略};
\node[card, minimum width=1.95cm] (c5d) at (2.18,7.615) {多源跨工况\\诊断问题建模};
\node[card, minimum width=1.95cm] (c5e) at (5.49,7.615) {逐源联合分布\\最优传输对齐};
\node[card, minimum width=1.95cm] (c5f) at (8.80,7.615) {类条件源域\\可靠性矩阵构建};
\node[card, minimum width=1.95cm] (c5g) at (12.11,7.615) {领域对抗特征\\与教师先验约束};
\node[card, minimum width=1.95cm] (c5h) at (15.42,7.615) {教师先验保护\\诊断结果融合};

\foreach \a/\b in {c5a/c5b,c5b/c5c,c5d/c5e,c5e/c5f,c5f/c5g,c5g/c5h}{
    \draw[flow] (\a.east) -- (\b.west);
}
\node[note, anchor=east] (toexp) at (8.52,6.815) {三大方法主线在TEP数据集上实验验证};

% =====================================================
% Chapter 6
% =====================================================
\node[block=cSix, minimum width=16.45cm, minimum height=3.25cm] (bg6) at (\cx,4.94) {};
\node[chaptermaintitle] at (\cx,6.19) {第六章\quad 实验设计与结果分析};
\node[card, minimum width=2.05cm] (c6a) at (4.05,5.46) {TEP 数据集\\6个生产模式};
\node[card, minimum width=2.05cm] (c6b) at (7.20,5.46) {统一划分协议\\目标标签仅评估};
\node[card, minimum width=2.05cm] (c6c) at (10.40,5.46) {$m_1,m_2,m_5$\\三模式组合};
\node[card, minimum width=2.05cm] (c6d) at (13.55,5.46) {$m_1$--$m_6$\\六模式组合};

\node[bigcard, minimum width=2.75cm] (c6e) at (4.05,3.85) {6个单源\\单目标域任务};
\node[bigcard, minimum width=2.75cm] (c6f) at (8.80,3.85) {3个两源\\单目标域任务};
\node[bigcard, minimum width=2.75cm] (c6g) at (13.55,3.85) {6个五源\\单目标域任务};

\foreach \a/\b in {c6a/c6b,c6b/c6c,c6c/c6d}{
    \draw[flow] (\a.east) -- (\b.west);
}
\coordinate (c6bus) at (10.40,4.62);
\draw[draw=black, line width=.50pt] (c6c.south) -- (c6bus) -- (4.05,4.62);
\draw[auxflow] (4.05,4.62) -- (c6e.north);
\draw[auxflow] (8.80,4.62) -- (c6f.north);
\draw[auxflow] (c6d.south) -- (c6g.north);

% 章节级路线：保留左右贴边走线，避免穿过内容密集区。
\draw[routebar] (bg2.south) -- (routeTop) -- (\routeL,18.82) -- (\routeL,8.815);
\draw[routeflow] (\routeL,16.82) -- (bg3.west);
\draw[routeflow] (\routeL,8.815) -- (bg5.west);

\draw[routebar] (bg4.east) -- (\routeR,12.82) -- (\routeR,6.815) -- (\cx,6.815);
\draw[routebar] (bg5.east) -- (\routeR,8.815);
\draw[routeflow] (\cx,6.815) -- (bg6.north);

% =====================================================
% Chapter 7
% =====================================================
\node[block=cSeven, minimum width=16.45cm, minimum height=1.62cm] (bg7) at (\cx,2.00) {};
\node[chaptermaintitle] at (\cx,2.48) {第七章\quad 总结与展望};
\node[card, minimum width=2.40cm] (c7a) at (3.40,1.69) {三条方法链\\研究总结};
\node[card, minimum width=2.40cm] (c7b) at (6.95,1.69) {方法作用边界\\与局限};
\node[card, minimum width=2.40cm] (c7c) at (10.50,1.69) {可扩展实验\\与消融验证};
\node[card, minimum width=2.40cm] (c7d) at (14.05,1.69) {工业部署\\后续方向};

\draw[auxflow] (bg6.south) -- (bg7.north);

\pgfresetboundingbox
\path[use as bounding box] (\xmin,\ymin) rectangle (\xmax,\ymax);

\end{tikzpicture}
